\documentclass{article}

% To pass options to natbib, set them before loading the DNU style:
% \PassOptionsToPackage{numbers,compress}{natbib}
\usepackage{dnu}

\usepackage[utf8]{inputenc}
\usepackage[T1]{fontenc}
\usepackage{hyperref}
\usepackage{url}
\usepackage{booktabs}
\usepackage{amsfonts}
\usepackage{nicefrac}
\usepackage{microtype}
\usepackage{xcolor}

\title{DNU Research Report Title}

\author{%
  Author Name\\
  Deep Network Understanding Lab\\
  Dickinson College\\
  \texttt{youremail@dickinson.edu}
  % Add coauthors using \And between author blocks.
}

\begin{document}

\maketitle

\begin{abstract}
  Summarize the research question, approach, principal results, and significance
  of the work in a concise paragraph.
\end{abstract}

\section{Introduction}

Introduce the problem, explain its importance, and summarize the report's
contributions.

\section{Related Work}

See the DNU Lab web pages for detailed instructions about this section.

\section{Methods (optional)}

Explain the data, materials, methods, or analytical approach in enough detail
to support interpretation and reproducibility.

\section{Results (optional)}

Present the principal findings, using figures and tables where appropriate.

\section{Discussion (optional)}

Interpret the findings, describe limitations, and identify implications or
future work.

\section{Conclusion}

Summarize the report's main conclusions.

\begin{ack}
  Optionally, acknowledge contributions, funding, or other support as appropriate.
\end{ack}

\section*{References}

Add the report's references here, using a consistent citation and bibliography
style.

\appendix

\section{Use Of Artificial Intelligence}

See the DNU Lab web pages for detailed instructions about this section.


\section{Supplementary material (optional)}

Add optional supplementary methods, proofs, or results here. 

\end{document}
